Es sei

\mavergleichskettedisp
{\vergleichskette
{M
}
{ =} { \R^2 \setminus \N_+
}
{ =} { \R^2 \setminus \{ (1,0) , (2,0), (3,0), \ldots  \}
}
{ } {
}
{ } {
}
}
{}{}{,}
es werden also aus dem $\R^2$ die auf der $x$-Achse platzierten positiven natürlichen Zahlen herausgenommen. Dies ist eine offene Teilmenge im $\R^2$ und damit eine differenzierbare Mannigfaltigkeit. Diese ist wegzusammenhängend, da man beispielsweise jeden Punkt aus $M$ mit
\mathl{y \geq 0}{} durch einen geraden Weg mit
\mathl{(0,1)}{} und jeden Punkt aus $M$ mit
\mathl{y \leq 0}{} durch einen geraden Weg mit
\mathl{(0,-1)}{} verbinden kann und diese beiden Punkte ebenfalls gerade verbindbar sind. Es sei nun

\mavergleichskettedisp
{\vergleichskette
{P
}
{ =} { (0,0)
}
{ } {
}
{ } {
}
{ } {
}
}
{}{}{}
und

\mavergleichskettedisp
{\vergleichskette
{M'
}
{ =} { M \setminus \{P\}
}
{ } {
}
{ } {
}
{ } {
}
}
{}{}{.}
Die Abbildung
\maabbeledisp {} {\R^2} {\R^2
} {(x,y)} {(x-1,y)
} {,}
ist
\zusatzklammer {als Verschiebung} {} {}
ein Diffeomorphismus und das Bild von $M$ ist genau $M'$. Daher sind
\mathkor {} {M} {und} {M'=M \setminus \{P\}} {}
zueinander diffeomorph.

 [[Kategorie:Latexseite]]
